\section*{Abstract}

Replace this paragraph with the abstract of your thesis. The abstract should briefly state the problem you addressed, the methods you used, the key results you obtained, and the conclusions you drew. Keep it concise — typically one paragraph, around 150 to 250 words.

This template is meant to provide an easy onboarding with \LaTeX.
We recommend using VSCode with the remote containers extensions.
Once you open this repository, you are asked to reopen in a container.

This container provides many useful features, such as
\begin{itemize}
    \item Automatic \LaTeX compilation on save.
    \item PDF Preview on the side. You can click in the PDF to find the \LaTeX source code, or you can sync our PDF view with your current Tex-File.
    \item Many useful keyboard shortcuts, e.g. for \textbf{bold} or \emph{italic} text.
    \item Automatic \textbf{grammar check} using Ltex Extension. It uses LanguageTool internally.
    \item Automatic \textbf{spell check}. You will see little yellow indicators in the editor.
    \item Many more\ldots you can provide pull requests to list your favorites features here.
    \item \emph{Note:} if you use VSCodium/Zed/\ldots instead of VSCode, some customization options like the extension need to be installed manually.
\end{itemize}

How can I improve my scientific writing skills?
\begin{itemize}
 \item \url{https://www.zfw.rub.de/sz/}
 \item \url{http://www.cs.joensuu.fi/pages/whamalai/sciwri/sciwri.pdf}
 \item \url{https://www.student.unsw.edu.au/writing}
 \item \url{https://www.sydney.edu.au/content/dam/students/documents/learning-resources/learning-centre/writing-a-thesis-proposal.pdf}
\end{itemize}